%% ma: L11-B15-0080
%% lop: 11
%% chuong: 5
%% bai: 15
%% dang: 11.15.5
%% loai: TLN
%% muc: VD
%% dap_an: "16"
%% nguon: "tai-lieu-hoc-tap-mon-toan-lop-11-theo-chuong-trinh-giao-duc-pho-thong-2018.pdf (Lớp Toán Thầy Duy, Phạm Lê Duy, THPT Chu Văn An), Chương 3 Bài 1, Phần 3, Câu 2 (tr.340)"
%% kien_thuc: "Lớp 11 Bài 15 (giới hạn của dãy số)"
%% trang_thai: cho-duyet
%% ly_do: "Dạng mới đề xuất (11.15.x), chờ giáo viên duyệt dạng; tài liệu gốc không có lời giải mẫu, đáp án do kho tự giải và kiểm bằng sympy."
\begin{ex}
Biết rằng giới hạn \[\lim\left(\sqrt{n^2+2n-2}-n\right)=\lim\dfrac{a\cdot\left(1-\dfrac1n\right)}{\sqrt{1+b\cdot\left(\dfrac1n-\dfrac1{n^2}\right)}+1}\] với $a;b$ là các số tự nhiên. Tính $S=(a+b)^2$.
\shortans{16}
\loigiai{\sloppy \[\sqrt{n^2+2n-2}-n=\dfrac{2n-2}{\sqrt{n^2+2n-2}+n}=\dfrac{2\left(1-\frac1n\right)}{\sqrt{1+2\left(\frac1n-\frac1{n^2}\right)}+1}.\] Vậy $a=2$, $b=2$ và $S=(2+2)^2=16$.}
\end{ex}
